% Transcription — fonds Alexandre Grothendieck, shelfmark 118, batch 1 (pages 1-20).
% Established from the facsimile archives/batches/118/batch-01.pdf (a mirror of
% https://grothendieck.umontpellier.fr/118.pdf), per the skill
% transcribe-grothendieck. The batch file carries no cover sheet: PDF page N
% is page N of the folder, and the archivists' pencil numbers bottom-left
% agree wherever legible. The folder is forty pages; this is the first of two
% batches.
%
% Pass: Opus 5.5 (claude-opus-5-5), 2026-09-24 — first pass, unchecked against the pages by a human.
%
% Thirteen of the twenty pages are transcribed: 2, 4, 7, 8, 10, 11, 13, 14,
% 15, 16, 18, 19, 20. Absent, with no \page and no note:
%   1      a pink cover in his hand, « Cohomology properties of maps in (Cat) :
%          1) Fibration-like properties » — no \page; his words become the
%          section and subsection headings, with a \note saying so;
%   3      an administrative accounts sheet (amounts, « Fonctionnement »),
%          nothing of his;
%   5, 6   ruled administrative tables (seminar names, amounts) in another
%          hand, nothing of his (page 7's paper shows the same table through);
%   9      a kraft envelope, printed « DU LANGUEDOC », addressed in pencil to
%          him — no mathematics;
%   12     a sheet inscribed « 1976 » and « Poèmes » in pencil, the first word
%          overwritten — unrelated to the mathematics;
%   17     a typed university notice of a doctoral defence, dated « Montpellier,
%          le 29 avril 1985 » — the only dated leaf in the batch; it dates
%          nothing that is transcribed and is recorded here only.
% No page is another person's text: nothing is summarised.
%
% What the batch is about. Working notes, mostly in English (pages 13, 16
% partly, 2, 4, 7, 10, 11) with French on pages 8, 14-15 and 18-20, on the
% « cohomological » classes of functors in (Cat) that the cover announces.
%   2      Three lattices without prose, the first two in ink, of classes of
%          maps: F, C (fibrations, cofibrations), Sm, Pr (smooth, proper),
%          S (Serre), their intersections, SprB, SprF, SprC; then the same
%          lattice intersected with W (weak equivalences), T = trivial,
%          Bias, As = W ∩ SprC, As° = W ∩ SprF; a third in pencil for
%          « strh » (strict homotopisms, cf. page 7).
%   4      The fullest version of these lattices, with B = F ∩ C, WSF, WSC,
%          WSB, Fl(Cat), and the second lattice with Tot As = W ∩ WSB; the
%          stability of each class under base change (proper, smooth, S, any)
%          and under intersection, with exceptions; « Sm ∩ Pr ⊂ S »,
%          « WSB ≠ WSF ∩ WSC ! »; « A Serre map is in W iff its fibers are
%          aspheric, i.e. S ∩ W = S ∩ T »; a question (*) asking for a
%          counterexample to stability of WSB under base change by T ∩ C.
%   7      Two words: « strict homotopism », by a struck « TF ».
%   8      On kraft paper, in French: characterisations of C/F, Pr/Sm (via
%          X_b → X_{/b} coaspherical, i.e. in As°, and commutation of R^i f_*
%          with base change), S (base change preserves W), WSC/WSF
%          (X_{/b} → X_{/b'} in W; Cofib(X/Y) → Y in S), WSB (3bis) via the
%          categories b_0\X/b_1; the inclusions F ⊊ Sm, C ⊊ Pr, Sm ∩ Pr ⊊ S
%          and the dualities C = F°, S = S°, Sm = Pr°. The leaf is torn off
%          at the foot.
%   10-11  Canonical factorisations of f : X → Y: Y ← Fib(f) ⇄ X (f^♭ a
%          fibration « but not yet Serre », μ_f mono, ρ_f a retraction), the
%          dual through Cofib(f); a cross of classes (Serre / Fibrations +
%          Cofibrations / monos / discrete, with smooth, proper, open and
%          closed immersions); the mapping cylinders Y → C(f) → X, C'(f) with
%          closed/open sections, C(Y), C'(Y), C'(f)° = C(f°). Page 11 is a
%          fair copy of page 10 adding the « trivial » row and the notes
%          discrete ∩ Serre = étale covering = smooth ∩ proper Serre, etc.
%   13     A table: f^♯ ∈ fibration ∩ mono, σ_f, λ_f, μ_f, ρ_f, f^♭ and their
%          classes; adjunctions (λ_f, σ_f)?, (μ_f, ρ_f); open and closed
%          immersions as maps to Δ_1; discrete (co)fibrations as presheaves.
%   14-15  A French programme after « l'histoire d'asphéricité » and Loc 1-3:
%          (0) Loc axioms, (1) duality, (2) commutation to fibrations,
%          cofibrations, immersions, cylinders, cones, (3) « sorite des ∫ »
%          of a span Y → X_0, Y → X_1 (functoriality, homotopisms), (4) a
%          theorem on amalgamated sums (cube), a question on pushouts along
%          open/closed immersions, the A^ version of the theorem
%          (i ∈ W_A ⇒ i' ∈ W_A, g ∈ W_A ⇒ f ∈ W_A), (5) Hot(W) and Hot_A,
%          (6) strong saturation of W, W_A.
%   16     « Sorite des recollements »: a category X glued from an open U and
%          closed Y through H_X : U° × Y → Ens; maps into Z, maps between
%          glued categories, products, homotopies compatible with the
%          decomposition.
%   18-20  His own run, numbered « 1 », « 2 », « 2 bis »: when a square
%          Y ⊂ X ⊃ U over Y' ⊂ X' ⊃ U' with U ≅ U' is cocartesian (no in
%          general; yes if f is cofibrant); gluing data for a functor
%          X' → Cat; ∫(Y → X, Y → Y') → X' = X ⊔_Y Y' should be in W, even
%          aspheric; it is a cofibration with fibres e or C(Y_{y'}); page 20
%          states the Theorem (cocartesian square in Cat, i', f cofibrations,
%          fibres equivalent to the point) and the Corollary i ∈ W ⇒ i' ∈ W,
%          g ∈ W ⇒ f ∈ W, via the fibres ∫(A = A, A → e), aspheric because
%          C(A) is a deformation retract of C(A) ⊔_A (A × Δ^1).
%
% Notation fixed for the batch. His class letters stay as written, in math
% italic: F, C, Sm, Pr, S, B, T, W, and the run-together WSF, WSC, WSB;
% \mathrm{Spr}, \mathrm{As}, \mathrm{Bias}, \mathrm{Tot}, \mathrm{strh},
% \mathrm{Fl}, \mathrm{Fib}, \mathrm{Cofib}, \mathrm{crib}, \mathrm{Cocrib}.
% On page 8 he writes the letters bold and repassed; \mathbf is kept there.
% X_{/b} and {}_{b\backslash}X for his over- and under-categories; \underline{W}
% for his underlined W (pages 14-20); (\mathrm{Cat}), (\mathrm{Ens}),
% (\mathrm{Sets}) with his parentheses; \Delta_1 / \Delta^1 as he places the
% index; e for the final object; A^{\wedge} for « A^ »; X^{\circ} for the
% opposite. The exponents of f read ♭ (f^♭, « Fib ») and ♮ (f^♮, the
% cylinder) and ♯ (page 13) are the least certain glyphs of the batch.
% « = » for « même » on page 19 is expanded with \add.
%
% Readings to check first: the class letter at the head of page 8, read C
% but also readable E (and « WSC » in rubric (3)); the letter read Pr in
% rubric (2) of page 8; « R^i f_* » / « R^i g_* » on page 8; the title of the
% « 2) » block on page 10; the exponents ♭, ♮, ♯; « UW » in W ∩ S = T ∩ S =
% UW (pages 2, 4); the heavily worked lower half of page 14 and the struck
% passage of page 20, both given largely as \ill.
%
% Cross-references. Page 11 recopies page 10; page 13's NB repeats page 11's.
% Page 2's « strh » is glossed by page 7. Batch 2 (pages 21-40) is transcribed
% separately.
\documentclass[11pt,a4paper]{article}
% Le préambule commun à toutes les transcriptions du fonds Grothendieck.
%
% Il tient en une idée : **l'incertitude se compose, elle ne se résout pas.**
% Une transcription qui lisse un mot douteux a détruit la seule chose pour
% laquelle on la faisait — un lecteur doit pouvoir savoir, à chaque endroit,
% ce qui était sur le feuillet et ce qui a été ajouté. Les six macros qui
% suivent sont donc l'appareil critique, et non de la décoration.
%
% Les mêmes commandes sont reconnues par `scripts/render.mjs`, qui produit la
% vue de lecture du site. Une macro ajoutée ici doit l'être aussi là-bas,
% faute de quoi le rendu échouera bruyamment — ce qui est voulu : un rendu
% silencieusement faux est pire qu'un rendu absent.

\ProvidesPackage{grothendieck}[2026/08/08 Transcriptions du fonds Grothendieck]

\usepackage{amsmath,amssymb,amsthm}
\usepackage{mathrsfs}
\usepackage{stmaryrd}
% Les diagrammes commutatifs sont partout dans ce fonds : c'est la langue
% même de la géométrie algébrique de Grothendieck, et une transcription qui
% les rendrait en prose ne transcrirait rien.
\usepackage{tikz-cd}
% Français par défaut — c'est la langue des feuillets — mais l'anglais est
% chargé aussi : la traduction et le résumé sont des documents anglais, et la
% typographie française (espaces avant les deux-points) y serait une faute.
% Ces documents basculent par \selectlanguage{english} en tête de corps.
\usepackage[english,french]{babel}
\usepackage{fontspec}
\usepackage{geometry}
\geometry{a4paper,margin=25mm,marginparwidth=28mm}
\usepackage{marginnote}
\usepackage{xcolor}
% ulem, not soul. `soul`'s \st and \ul are built on a character-by-character
% scan that XeLaTeX's font machinery breaks: the document fails with an
% undefined control sequence rather than degrading. ulem does the same two jobs
% and is Unicode-engine safe. `normalem` stops it hijacking \emph, which the
% transcriptions use for Grothendieck's own emphasis.
\usepackage[normalem]{ulem}
\usepackage{hyperref}
\hypersetup{colorlinks=true,linkcolor=black,urlcolor=black,pdfborder={0 0 0}}

% --- Métadonnées du lot -----------------------------------------------------
% Reprises telles quelles par le script de rendu, qui les affiche en tête de
% la vue de lecture. Elles fixent ce que le fichier prétend couvrir : sans
% elles, une transcription flotte sans point d'attache dans un fonds de seize
% mille feuillets.

\newcommand{\folder}[1]{\gdef\grothFolder{#1}}
\newcommand{\batch}[1]{\gdef\grothBatch{#1}}
\newcommand{\leaves}[2]{\gdef\grothFirst{#1}\gdef\grothLast{#2}}
\newcommand{\foldertitle}[1]{\gdef\grothFoldertitle{#1}}
\gdef\grothFolder{?}\gdef\grothBatch{?}\gdef\grothFirst{?}\gdef\grothLast{?}\gdef\grothFoldertitle{}

% --- Le résumé --------------------------------------------------------------
%
% Il ouvre la lecture modernisée, sous le titre « Résumé », et il est composé
% en petit. Ce n'est pas un ornement : c'est le seul passage du document qui
% ne soit pas des mathématiques, et un lecteur doit pouvoir le reconnaître
% avant de l'avoir lu — puis le sauter s'il connaît déjà le sujet, ou s'y
% arrêter s'il ne le connaît pas. Le filet qui le suit marque la reprise du
% corps.

\newenvironment{resume}
  {\small\setlength{\parskip}{0.45em}}
  {\par\medskip\hrule\medskip}

% --- Mots-clés ---------------------------------------------------------------
%
% En anglais, en clôture du résumé de la lecture modernisée : le vocabulaire
% moderne sous lequel chercher ce que les feuillets construisent. Le manifeste
% du site (`npm run manifest`) les extrait pour étiqueter la cote entière —
% cette ligne est la source unique des tags, il n'y a pas de fichier à côté.
\newcommand{\keywords}[1]{\par\smallskip\noindent
  {\footnotesize\textsc{Keywords} --- \textit{#1}}\par}

% --- L'appareil critique ----------------------------------------------------

% Le numéro de feuillet, en marge. C'est l'ancre qui permet de régler un
% désaccord sur une formule en regardant *un* feuillet plutôt qu'une chemise
% de six cents. Le site s'en sert aussi pour tourner les pages du fac-similé
% quand on fait défiler la transcription.
\newcommand{\leaf}[1]{%
  \marginnote{\footnotesize\color{gray!70}\textsf{#1}}%
  \label{leaf:#1}\ignorespaces}

% Illisible : on ne devine pas. Un mot inventé qui se lit comme les autres est
% le pire résultat possible.
\newcommand{\ill}{\textcolor{red!60!black}{[\,\ldots\,]}}

% Lecture douteuse : le mot est donné, et signalé comme incertain.
\newcommand{\uncertain}[1]{\uline{#1}}

% Ajout de l'éditeur — une lettre manquante, un mot sous-entendu.
\newcommand{\add}[1]{\textcolor{blue!50!black}{[#1]}}

% Biffé par Grothendieck lui-même. Ce qu'il a rayé fait partie du document :
% une rature dit souvent où la pensée a changé de direction.
\newcommand{\struck}[1]{\sout{#1}}

% Note du transcripteur, sur l'état matériel du feuillet.
\newcommand{\note}[1]{\par\smallskip{\small\color{gray!60!black}\textit{[#1]}}\par\smallskip}

% Note marginale de Grothendieck, distincte du corps du texte.
\newcommand{\marginal}[1]{\par\smallskip\noindent\hspace{1em}%
  {\small\color{gray!60!black}#1}\par\smallskip}

% --- Titre ------------------------------------------------------------------

\AtBeginDocument{%
  \begin{center}
    {\large\bfseries Fonds Alexandre Grothendieck --- cote n\textsuperscript{o}~\grothFolder}\\[2pt]
    {\itshape\grothFoldertitle}\\[4pt]
    {\small feuillets \grothFirst--\grothLast\ (lot \grothBatch)}\\[2pt]
    {\footnotesize\color{gray!60!black}Transcription établie d'après le fac-similé numérisé
     par l'Université de Montpellier.\\
     Les lectures incertaines sont soulignées, les illisibles notées [\,\ldots\,],
     les ajouts entre crochets.}
  \end{center}
  \vspace{1em}}

\endinput


\folder{118}
\batch{1}
\pages{1}{20}
\dating{[à partir de 1983]}
\watermark{Édition de démonstration}
\foldertitle{Cohomology properties of maps in (Cat) : notes manuscrites (s.d.).}

\begin{document}

\section{Cohomology properties of maps in (Cat)}
\note{titre de sa main, sur la chemise rose qui forme la page 1, suivi de deux points ; la même chemise porte la ligne suivante, reprise en sous-titre}

\subsection{1) Fibration-like properties}
\note{« like » est lu sans certitude}

\page{2}
\note{page sans texte : trois treillis de classes de morphismes de $(\mathrm{Cat})$, les deux premiers à l'encre, le troisième au crayon. Les lettres sont les siennes : $F$, $C$ (fibrations, cofibrations), $Sm$, $Pr$ (lisse, propre), $S$ (sans doute Serre), les primes $F'$, $C'$ étiquetant les flèches ; $\mathrm{Spr}$ est écrit tel quel}

\begin{tikzcd}[column sep=tiny, row sep=small, nodes={font=\scriptsize}]
 & & & & F\cap C \arrow[dl, "F\cap C'"'] \arrow[dr, "F'\cap C"] & & & & \\
 & & & F\cap Pr \arrow[dl] \arrow[dr, "F'\cap Pr"] & F'\cap C' & Sm\cap C \arrow[dl, "Sm\cap C'"'] \arrow[dr] & & & \\
 & & F\cap S \arrow[dl] \arrow[dr, "F'\cap S"] & & Sm\cap Pr \arrow[dl] \arrow[dr] & & S\cap C \arrow[dl, "S\cap C'"'] \arrow[dr] & & \\
 & F \arrow[dr, "F'"] & & Sm\cap S \arrow[dl] \arrow[dr] \arrow[dddd, dashed] & & S\cap Pr \arrow[dl] \arrow[dr] \arrow[dddd, dashed] & & C \arrow[dl, "C'"] & \\
 & & Sm & & S \arrow[d] & & Pr & & \\
 & & & & \mathrm{Spr}B \arrow[d] & & & & \\
 & & & & \mathrm{Spr}F\cap \mathrm{Spr}C \arrow[dl] \arrow[dr] & & & & \\
 & & & \mathrm{Spr}F & & \mathrm{Spr}C & & &
\end{tikzcd}

\note{$F'\cap C'$, au centre du premier losange, n'est l'étiquette d'aucune flèche ; les deux flèches en pointillé descendent de $Sm\cap S$ vers $\mathrm{Spr}F$ et de $S\cap Pr$ vers $\mathrm{Spr}C$}

\begin{tikzcd}[column sep=0pt, row sep=small, nodes={font=\tiny}]
 & & W\cap F\cap C \arrow[dl] \arrow[dr] & & \\
 & W\cap F\cap Pr \arrow[dl] \arrow[dr] & & W\cap Sm\cap C \arrow[dl] \arrow[dr] & \\
 W\cap F\cap S = T\cap F \arrow[dr] & & W\cap Sm\cap Pr \arrow[dl] \arrow[dr] & & W\cap S\cap C = T\cap C \arrow[dl] \\
 & W\cap Sm\cap S = T\cap Sm \arrow[dr] & & W\cap S\cap Pr = T\cap Pr \arrow[dl] & \\
 & & W\cap S = T\cap S = \mathcal{U}W \arrow[d] & & \\
 & & \mathrm{Bias} = W\cap \mathrm{Spr}B \arrow[d] & & \\
 & & \mathrm{As}\cap \mathrm{As}^{\circ} \arrow[dl] \arrow[dr] & & \\
 & \mathrm{As}^{\circ} = W\cap \mathrm{Spr}F \arrow[dr] & & \mathrm{As} = W\cap \mathrm{Spr}C \arrow[dl] & \\
 & & W & &
\end{tikzcd}

\note{le signe lu $\mathcal{U}W$ (dans $T\cap S = \mathcal{U}W$) est incertain}

\begin{tikzcd}[column sep=tiny, row sep=small, nodes={font=\scriptsize}]
 & & \mathrm{strh}\cap F\cap C \arrow[dl, no head] \arrow[dr, no head] & & \\
 & \mathrm{strh}\cap F\cap \mathrm{str}Pr \arrow[dl, no head] \arrow[dr, no head] & & \mathrm{strh}\cap \mathrm{str}Sm\cap C \arrow[dl, no head] \arrow[dr, no head] & \\
 \mathrm{strh}\cap F \arrow[dr, no head] & & \mathrm{strh}\cap \mathrm{str}Sm\cap \mathrm{str}Pr \arrow[dl, no head] \arrow[dr, no head] & & \mathrm{strh}\cap C \arrow[dl, no head] \\
 & \mathrm{strh}\cap \mathrm{str}Sm \arrow[dr, no head] & & \mathrm{strh}\cap \mathrm{str}Pr \arrow[dl, no head] & \\
 & & \mathrm{strh} & &
\end{tikzcd}

\note{ce troisième treillis est au crayon, ses traits sans pointes ; « strh » est sans doute « strict homotopism » (voir la page 7). Un long trait au crayon relie $\mathrm{strh}$ au bas du deuxième treillis, sans qu'on voie à quel sommet}

\page{4}
\note{feuille de travail en anglais, encre et crayon mêlés : deux treillis de classes de morphismes au centre, des remarques encadrées ou cernées tout autour. On donne d'abord les treillis, puis les remarques, de haut en bas et de gauche à droite. Les notations $F$, $C$, $Sm$, $Pr$, $S$ sont celles de la page 2 ; $WSF$, $WSC$, $WSB$ sont écrits d'un seul tenant. Des lignes en tirets, au crayon, doublent chacun des deux treillis et ne sont pas reproduites}

\marginal{$S$, $Pr$, $Sm$ : Serre map i.e.\ Hot-fibration, proper, smooth}
\note{en haut, au crayon, chacune des trois lettres reliée par un trait à son nom}

\marginal{$T =$ « universally in $W$ »}
\note{en haut à droite}

\begin{tikzcd}[column sep=tiny, row sep=small, nodes={font=\scriptsize}]
 & & & & B \overset{\mathrm{def}}{=} F\cap C \arrow[dl] \arrow[dr] & & & & \\
 & & & F\cap Pr \arrow[dl] \arrow[dr] & & Sm\cap C \arrow[dl] \arrow[dr] & & & \\
 & & F\cap S \arrow[dl, "F\cap WSC"'] \arrow[dr] & & Sm\cap Pr \arrow[dl] \arrow[dr] & & S\cap C \arrow[dl] \arrow[dr, "WSF\cap C"] & & \\
 & (F) \arrow[dr] & & Sm\cap S \arrow[dl, "Sm\cap WSC"'] \arrow[dr] \arrow[dddd] & & S\cap Pr \arrow[dl] \arrow[dr, "WSF\cap Pr"] \arrow[dddd] & & (C) \arrow[dl] & \\
 & & (Sm) & & (S) \arrow[d] & & (Pr) & & \\
 & & & & \boxed{WSB} \arrow[d] & & & & \\
 & & & & WSF\cap WSC \arrow[dl] \arrow[dr] & & & & \\
 & & & WSF \arrow[dr] & & WSC \arrow[dl] & & & \\
 & & & & \mathrm{Fl}(\mathrm{Cat}) & & & &
\end{tikzcd}

\note{les parenthèses rendent les cercles dont il entoure $F$, $C$, $Sm$, $S$, $Pr$. Les deux longues flèches verticales, de $Sm\cap S$ vers $WSF$ et de $S\cap Pr$ vers $WSC$, sont hachurées de petits traits. Sur la flèche $F\cap S \to F$, le mot « str » ; les étiquettes $F\cap WSC$, $Sm\cap WSC$, $WSF\cap C$, $WSF\cap Pr$ sont au crayon, chacune attachée par un point à sa flèche. Au-dessus de $WSF\cap WSC$, une première inscription est noircie et ne se lit pas}

Sous $WSF$ : (\struck{\ill{}} Serre \emph{fibrations}), stable by proper b.ch. Sous $WSC$ : $WSC$ or (\struck{\ill{}} Serre precofibrations), stable by smooth b.ch.
\note{« precofibrations » est une lecture incertaine ; $WSF$ et $WSC$ sont repassés sur d'autres lettres}

\marginal{stable under any base change ; squares cartesian}
\note{à droite du premier treillis, reliée à lui par une accolade}

\marginal{\ill{} non ch. \quad stable by $\Delta_1$ base ch.}
\note{à droite, en oblique, sous deux traits en accolade}

\begin{tikzcd}[column sep=0pt, row sep=small, nodes={font=\tiny}]
 & & \mathrm{str}\ W\cap F\cap C \arrow[dl] \arrow[dr] & & \\
 & \mathrm{str}\ W\cap F\cap Pr \arrow[dl] \arrow[dr] & & W\cap Sm\cap C \arrow[dl] \arrow[dr] & \\
 W\cap F\cap S = T\cap F \arrow[dr] & & W\cap Sm\cap Pr \arrow[dl] \arrow[dr] & & W\cap S\cap C = T\cap C \arrow[dl] \\
 & W\cap Sm\cap S = T\cap Sm \arrow[dr] & & W\cap S\cap Pr = T\cap Pr \arrow[dl] & \\
 & & W\cap S = T\cap S = UW \arrow[d, "\mathrm{str}"] & & \\
 & & \mathrm{Tot}\,\mathrm{As} = W\cap WSB \arrow[dl, "\mathrm{str}"'] \arrow[dr] & & \\
 & \mathrm{As}^{\circ} = W\cap WSF \arrow[dr, "\mathrm{str}"'] & & \mathrm{As} = W\cap WSC \arrow[dl] & \\
 & & W & &
\end{tikzcd}

\note{$W\cap S$ et $\mathrm{As}$ sont cerclés deux fois, $\mathrm{As}^{\circ}$ une fois, $\mathrm{Tot}\,\mathrm{As}$ en pointillé ; sur $\mathrm{As}^{\circ}$ une première inscription, « str » suivi d'un mot, est noircie ; devant $= T\cap F$, un signe biffé}

Sous $\mathrm{As}^{\circ}$ : (stable by proper b.ch.) ; sous $\mathrm{As}$ : (stable by smooth b.ch.) ; sous $W$ : (stable by $S$ b.ch.).

\marginal{stable under any base change ; squares cartesian}
\note{à droite du second treillis, reliée à lui par une accolade}

\marginal{stable by intersection except $WSF\cap C$ and $WSC\cap Pr$, and dually $F\cap WSC$, $Sm\cap WSC$ and $WSF\cap WSC$}
\note{entre les deux treillis, à droite, au crayon puis repassé à l'encre}

\marginal{stable by intersection except $\mathrm{As}\cap \mathrm{As}^{\circ}$}
\note{en bas à droite}

Encadrés, en haut à gauche :
\[
  \boxed{Sm\cap Pr \subset S} \qquad
  \boxed{\begin{array}{l} WSC\cap Pr \subset S \\ WSF\cap Sm \subset S \end{array}}
\]
(hence $WSC\cap Pr = (S\cap Pr)$) \quad ( — $WSF\cap Sm = (S\cap Sm)$ ) \quad $(WSB\cap Sm)$
\note{dans la seconde égalité, le $S$ est repassé sur un autre mot ; « $(WSB\cap Sm)$ » est écrit dessous. En tête de page, au-dessus, « $WSB\cap$ \struck{\ill{}} $=$ », relié par une flèche courbe à « $(S\cap Pr$ »}

\struck{$WSB\cap Sm \not\subset S$ ? \quad ($WSB\cap Pr \not\subset WSC$)}

$WSB \neq WSF\cap WSC$ !

$WSC$, $WSB$ not stable under $S\cap C$ base ch.? (And dually!) even if both are « trivial »?$^{*}$
\note{ces trois dernières lignes sont enfermées dans un même trait}

(\uncertain{$B$} $\not\subset WSF, WSC$ ! \quad \struck{\ill{}} $[C \not\subset WSF,\ F \not\subset WSC]$

$C\cap WSF \not\subset S$ ? \quad $F\cap WSC \not\subset S$ ?
\note{la première lettre, lue $B$, est surchargée}

A Serre map is in $W$ iff \struck{\ill{}} its fibers are aspheric, i.e.\ $S\cap W = S\cap T$, and also iff it is univ.\ in $W$.

Even if fibers are aspheric, the Serre condition cannot be tested by only taking base-changes $\Delta_1 \to X$ ?

$^{*}$ i.e.\ find $f : X \to Y$, $g : Y' \to Y$ with $f \in W\cap WSB$ (i.e.\ $f$ tot.\ aspher.) \struck{$W$} and $g \in T\cap C$ (cofibration with trivial fibers) s.th.\ $f' : X' \to Y'$ is not in \uncertain{$WSB$} (say $Y'$ with final element and $y' \in \mathrm{Ob}\, Y'$ s.th.\ $X'_{/y'} \to X'_{/e}$ not a weak equiv.?)
\note{les trois derniers paragraphes sont chacun cernés d'un trait ; le dernier renvoie à l'astérisque de « trivial »}

this square non cartesian? $\longrightarrow$
\note{cerné, en bas à gauche ; la flèche vise le carré inférieur du second treillis}

\page{7}
\note{au crayon, au milieu d'une page par ailleurs blanche}

$TF$ \struck{\ill{}} \quad $X \downarrow Y$ \quad \emph{strict homotopism}
\note{$X \downarrow Y$ est une flèche verticale de $X$ vers $Y$, écrite sous $TF$ ; « strict homotopism » est souligné d'un trait ondulé}

\page{8}
\note{sur papier kraft, à l'encre, en français cette fois ; la page est divisée par des traits horizontaux en rubriques numérotées et cerclées, chacune coupée par un trait vertical : à gauche une classe de morphismes, à droite la classe duale. Pour $f : X \to Y$ et $b$ objet de $Y$, il écrit $X_{/b}$ et ${}_{b\backslash}X$ ; on garde ces graphies. Les équivalences $\Leftrightarrow$ qu'il enchaîne verticalement sont rendues par « $\Leftrightarrow$ » en tête de ligne}

$f : X \to Y$

$f \in \mathbf{C} \Longleftrightarrow X_{/b} \to X_{/b'}$ cofinal à gauche
\note{le numéro cerclé, en marge, est emporté par la déchirure. La lettre, repassée, se lit aussi bien $E$ ; on lit $C$ parce que la colonne de droite porte $F$, comme en (3) $WSC$ fait face à $WSF$ — lecture à vérifier, de même en (3)}

$\Leftrightarrow$ ${}_{b'\backslash}X \to {}_{b\backslash}X$ à fibres non vides \uncertain{ayant} un objet initial

$\Leftrightarrow$ La catégorie des relèv.\ d'une flèche $b \to b'$, de source donnée $a \in X_b$, a objet initial

$f \in F \Longleftrightarrow$ \quad $\Updownarrow$ \quad $\Updownarrow$
\note{colonne de droite laissée vide, sauf ces deux signes}

(2) $f \in Pr \Longleftrightarrow (X_{b} \to X_{/b})$ strongly \uncertain{équivalence} coasph., i.e.\ $\in \mathrm{As}^{\circ}$
\note{la lettre après $\in$ est surchargée (un $F$ sur $Pr$, ou l'inverse) ; on lit $Pr$, la colonne de droite portant $Sm$. Le premier $X$ est repassé ; « strongly » est lu sans certitude, comme le mot suivant}

$\Leftrightarrow$ ${}_{a\backslash}X \to {}_{b\backslash}Y$ à fibres asph.

$\Leftrightarrow$ \struck{\ill{}} $({}_{a\backslash}X \to {}_{b\backslash}Y) \in \mathbf{S}$ \struck{\ill{}}

$\Leftrightarrow$ La cat.\ des rel.\ d'une flèche $b \to b'$, de source donnée $a \in \mathrm{Ob}\, X_b$, est asph.

$\Leftrightarrow$ \uncertain{$R^i f_{*}$} commute aux chgt de base quelc.

\marginal{${}_{a\backslash}X \to {}_{b\backslash}Y$ is proper (with \uncertain{discrete} fibres)}
\note{cerné, en marge gauche, une flèche le rattachant à la troisième ligne}

$f \in Sm \Longleftrightarrow$ \quad $\Updownarrow$ \quad $\Updownarrow$ \quad le chgt de base par $f$ commute à \uncertain{$R^i g_{*}$}
\note{colonne de droite. Les $R^i f_{*}$, $R^i g_{*}$ de la rubrique (2) sont des lectures incertaines d'abréviations serrées (« Rif », « Rig »)}

[ $\Longleftarrow$ ? base changes by $f$ preserves \struck{\ill{}} \uncertain{cocartesian} squares ] \quad ( $\Longleftarrow$ ?
\note{en anglais, au pied de la rubrique (2), sous toute sa largeur}

(1) $f \in \mathbf{S} \Longleftrightarrow \forall g : Y'' \to Y'$ sur $Y$, si $g \in W$, alors $(g \times_Y X : X'' \to X') \in W$.

\[
  \boxed{\begin{array}{l} F \subsetneq Sm \\ C \subsetneq Pr \end{array}} \qquad
  \boxed{Sm \cap Pr \subsetneq S} \qquad
  \left\{\begin{array}{l} C = F^{\circ} \\ F = C^{\circ} \end{array}\right. \quad
  S = S^{\circ} \qquad
  \left\{\begin{array}{l} Sm = Pr^{\circ} \\ Pr = Sm^{\circ} \end{array}\right.
\]
\[
  \boxed{T\cap Pr \subset S,\ T\cap Sm \subset S}
\]
\note{après $C = F^{\circ}$, un mot biffé}

(3) $f \in WS\mathbf{C} \Longleftrightarrow \forall b \to b'$ dans $Y$, $(X_{/b} \to X_{/b'}) \in W$ — (a)

$\Leftrightarrow$ $\forall Y'' \xrightarrow{\ u\ } Y'$ sur $Y$, $Y''$ et $Y'$ \struck{\ill{}} fibrés sur $W$, si $u \in W$, alors $(g \times_Y X : X'' \to X') \in W$ — (b)
\note{« fibrés » est écrit au-dessus d'un mot biffé et lu sans certitude}

$\Leftrightarrow$ \struck{\ill{}} $(\mathrm{Cofib}(X/Y) \to Y) \in$ \struck{\ill{}} $\mathbf{S}$ — (a')

stable par chgt de base par \struck{fibrations} \uncertain{propres lisses} (mais non par $g \in S$ ?)
\note{la ligne est soulignée deux fois ; les deux derniers mots sont écrits au-dessus du mot biffé}

\marginal{$R^i f_{*}$ (c\textsuperscript{t}) loc.\ const.}
\note{en marge gauche, à hauteur de (a'), avec une flèche ; lecture incertaine}

$f \in WSF \Longleftrightarrow$ \quad $\Updownarrow$ \quad $\Longleftrightarrow (\mathrm{Fib}(X/Y) \to Y) \in$ \uncertain{$S$} \quad stable par chgt de base \emph{propres}
\note{colonne de droite de (3)}

\[
  \boxed{WSC \cap Pr = S} \qquad WSF \cap Sm \subset S
\]
\note{le premier encadré se lit $WSC\cap Pr = S$, pour $\subset S$ sans doute, comparer la page 4}

(3bis) $f \in WS\mathbf{B} \Longleftrightarrow$ \struck{\ill{}} pour $b'_0 \to b_0 \xrightarrow{\ u\ } b_1 \to b'_1$ dans $Y$,
\note{« 3bis » cerclé ; $WS\mathbf{B}$ est repassé sur une autre lettre ; au-dessus, « $\Longleftarrow WSE$ » biffé. Sous $b'_0 \to b_0 \to b_1 \to b'_1$ une accolade porte $u'$, pour la flèche composée $b'_0 \to b'_1$ sans doute}

(def) \quad ${}_{b_0\backslash}X_{/b_1} \longrightarrow {}_{b'_0\backslash}X_{/b'_1}$ est $\in W$, i.e.\ l'équiv.\ faible ${}_{b_0\backslash}Y_{/b_1} \to {}_{b'_0\backslash}Y_{/b'_1}$ \struck{\ill{}} est transformée en équiv.\ faible par $X_{\times}(?)$.
\note{les deux catégories ${}_{b_0\backslash}X_{/b_1}$, ${}_{b'_0\backslash}X_{/b'_1}$ sont cernées, et reliées par des traits à « $X[u]$ » et « $X[u']$ », écrits à gauche, et au mot « fibre » entre elles}

$\Leftrightarrow$ $\forall b \to b'$ dans $Y$, ${}_{b'\backslash}X \to {}_{b\backslash}X$ est $\in W$ \struck{\ill{}} sur $Y$.

\struck{et} $\forall b \to b'$ dans $Y$, $X_{/b} \to X_{/b'}$ est $\in W$ \uncertain{coloc.} sur $Y$
\note{la page s'arrête là, au bord déchiré}

\page{10}
\note{sur papier rose, en anglais, au crayon et à l'encre}

\marginal{\struck{Coho.\ propr.\ of \ill{}} \quad 2) \uncertain{Immersions fermées et ouvertes}, \uncertain{propriétés cohomologiques} \ill{}}
\note{en haut à droite, dans un cadre : une première ligne rayée, puis ce titre de rubrique, qui fait pendant au « 1) » de la page 1 ; lecture incertaine}

\begin{tikzcd}[column sep=large]
  Y & \mathrm{Fib}(f) \arrow[l, "f^{\flat}"'] & X \arrow[l, hook', "\mu_f"'] \\
  & \mathrm{Fib}(f) \arrow[r, bend right=20, "\rho_f"'] & X
\end{tikzcd}

\note{rendu en deux lignes pour séparer les deux flèches entre $\mathrm{Fib}(f)$ et $X$ : sur la page elles sont superposées, $\mu_f$ (« mono ») de $X$ vers $\mathrm{Fib}(f)$, et au-dessous, en arc, $\rho_f$ de $\mathrm{Fib}(f)$ vers $X$}

$f^{\flat}$ : fibration (Fibers ${}_{b\backslash}X$ discretely cofibered$/X$) \emph{but not yet Serre}. \quad $\rho_f$ : (\struck{\ill{}} Serre fibering) retraction [fibers $Y_{/b}$ discr.\ fibered over $Y$].

\marginal{not cofib.\ nor fib, i.e.\ not cl.\ imm, nor open imm}
\note{au-dessus de $\mu_f$, avec une flèche vers « mono »}

\begin{tikzcd}[column sep=large]
  Y & \mathrm{Cofib}(f) \arrow[l, "f^{\flat\prime}"'] & X \arrow[l, hook', "\mu'_f"'] \\
  & \mathrm{Cofib}(f) \arrow[r, bend right=20, "\rho'_f"'] & X
\end{tikzcd}

$f^{\flat\prime}$ : cofibration (fibers $X_{/b}$ discretely fibered$/X$) \emph{but not yet Serre}. \quad $\rho'_f$ : Serre fibering retraction (fibers ${}_{b\backslash}Y$ discr.\ cofibered over $Y$) but not \ill{}.
\note{« mono » est écrit aussi au-dessus de $\mu'_f$ ; la dernière phrase s'arrête en bord de page}

\note{au milieu de la page, une croix à quatre bras, cadres tracés au crayon ; on la rend bras par bras}
\[
  \begin{array}{ccccc}
    & & \text{bifibration} & & \\
    & & \text{Serre} & & \\
    \text{Serre fibration} & \textbf{Fibrations} & + & \textbf{Cofibrations} & \text{Serre cofibration} \\
    & & \text{monoe} & & \\
    \text{open imm.} & & & & \text{closed immersion}
  \end{array}
\]
À gauche, dans un crochet : smooth Serre, smooth ; une flèche va de « smooth » à « Fibrations ». À droite, dans un crochet : proper Serre, proper ; une flèche va de « proper » à « Cofibrations ».
\note{« monoe » (sic) sur la page, pour « monos » sans doute}

\[
  Y \xrightarrow[\text{imm}]{\ f^{\natural}\ \text{open}\ } C(f) \xrightarrow{\ \lambda_f\ (\text{epi})\ } X,
  \qquad \sigma_f : X \to C(f) \ \text{closed section}
\]
\[
  Y \xrightarrow[\text{imm}]{\ f^{\natural\prime}\ \text{closed}\ } C'(f) \xrightarrow{\ \lambda'_f\ } X,
  \qquad \sigma'_f : X \to C'(f) \ \text{open section}
\]
\note{les deux sections sont dessinées en arc de retour sous la flèche $\lambda$ ; l'exposant de $f$, lu $\natural$, est incertain ; la flèche de $Y$ vers $C(f)$ porte un crochet d'inclusion}

\[
  Y \xrightarrow[\text{open imm}]{\ i_Y\ } C(Y) \xrightarrow{\ \varepsilon_Y\ } e
  \quad \text{closed section } (e \text{ strict final object})
\]
\[
  Y \xrightarrow[\text{closed imm}]{\ i'_Y\ } C'(Y) \xrightarrow{\ \varepsilon'_Y\ } e
  \quad \text{open section } (e \text{ strict initial object})
\]
\note{« strict » est ajouté au-dessus de « final » ; $i'_Y$ porte un point au-dessus, peut-être un accent}

\[
  C'(f)^{\circ} = C(f^{\circ}) \quad \text{i.e.} \quad C'(f) = C(f^{\circ})^{\circ}
\]
\[
  C'(Y)^{\circ} = C(Y^{\circ}) \quad \text{i.e.} \quad C'(Y) = C(Y^{\circ})^{\circ}
\]

\page{11}
\note{au crayon, en anglais ; reprise au propre de la page 10, avec les compléments « trivial … » et « homotopism »}

\begin{tikzcd}
  Y & \mathrm{Fl}(Y) \arrow[l] \arrow[d, bend right=15] & \mathrm{Fib}(f) \arrow[l] \arrow[ll, bend right=30, "f^{\flat}"'] \arrow[d, "\rho_f"] \\
  & Y' \arrow[u, dashed, bend right=15] & X \arrow[l, "f"] \arrow[u, dashed, bend right=40, "\mu_f"']
\end{tikzcd}

\note{en haut à gauche, petit carré au crayon ; les deux flèches entre $\mathrm{Fl}(Y)$ et $Y'$ sont l'une pleine vers le bas, l'autre en pointillé vers le haut, et $\mu_f$ est un arc en pointillé de $X$ vers $\mathrm{Fib}(f)$. Le $Y'$ se lit peut-être $Y$ ; l'étiquette $f$ semble barrée}

\begin{tikzcd}[column sep=large]
  Y & \mathrm{Fib}(f) \arrow[l, "f^{\flat}"'] \arrow[r, bend right=25, "\rho_f"'] & X \arrow[l, hook', "\mu_f\ \mathrm{mono}"']
\end{tikzcd}

$f^{\flat}$ : fibration (fibers ${}_{b\backslash}X$ discretely cofibered over $X$) \emph{but not yet Serre}. \quad $\rho_f$ : trivial Serre cofibering retraction (fibers $Y_{/b}$ discr.\ fibered over $Y$) and homotopism.

\note{au centre, une croix à quatre bras cernée d'un trait, et au-dessus un tableau à deux lignes ; on les rend ligne par ligne}

\[
  \scriptstyle
  \begin{array}{c|c|c|c|c}
    \text{smooth Serre} & \text{Serre fibration} & \text{Serre} & \text{Serre cofibration} & \text{Serre proper} \\
    \hline
    \text{trivial smooth} & \text{trivial Serre fibration} & \left[\begin{array}{c}\text{Serre}\cap W \\ = \text{triv Serre}\end{array}\right] & \text{trivial cofibration} & \text{trivial proper}
  \end{array}
\]
\note{sous « Serre », une flèche verticale descend vers $[\text{Serre}\cap W = \text{triv Serre}]$}

\[
  \begin{array}{ccccc}
    \text{smooth} & \hookleftarrow & \textit{fibration} \quad + \quad \textit{cofibration} & \hookrightarrow & \text{proper} \\
    & & \text{mono} & & \\
    & & \downarrow & & \\
    & & \text{discrete} & &
  \end{array}
\]
\note{sous « mono », un mot biffé et illisible}
\note{les flèches entre « smooth » et « fibration », « cofibration » et « proper » sont des inclusions ; leur sens est lu de la classe la plus petite vers la plus grande}

Sous le bras gauche :
\[
  \left.\begin{array}{l}
    \text{open immersion} \longleftrightarrow X^{\circ} \xrightarrow{\ \varphi\ } \Delta^1 = [0 \to 1] \\
    \phantom{\text{open immersion}} \longleftrightarrow \mathrm{crib}(X) \\
    \phantom{\text{open immersion \ }} Y = \varphi^{-1}(\{1\}) = \psi^{-1}(\{0\}) = X_0
  \end{array}\right\}
\]
(où $\psi : X \xrightarrow{\ \varphi^{\circ}\ } \Delta_1^{\circ} \overset{\mathrm{can}}{\simeq} \Delta_1$)

discrete fibration $\longleftrightarrow \hat{X}$, i.e.\ $X^{\circ} \to (\mathrm{Sets})$

Sous le bras droit :
\[
  \left\{\begin{array}{l}
    \text{closed immersion} \longleftrightarrow X \xrightarrow{\ \varphi\ } \Delta^1 = [\emptyset \to \{\emptyset\}] \\
    \phantom{\text{closed immersion}} \longleftrightarrow \mathrm{Cocrib}(X) \\
    \phantom{\text{closed immersion \ }} Y = \varphi^{-1}(\{1\}) = X_1
  \end{array}\right.
\]
\note{au-dessus de $\emptyset$ et $\{\emptyset\}$, les chiffres $0$ et $1$ ; sous « closed immersion », un signe $=$ isolé}

discrete cofibration $\longleftrightarrow \check{X}$, i.e.\ $X \to (\mathrm{Sets})$

\[
  Y \xrightarrow[\text{open imm}]{\ f^{\natural}\ } C(f) \xrightarrow{\ \lambda_f\ \text{epi}\ } X,
  \qquad \sigma_f : \text{closed section and \textit{homotopism}}
\]
\[
  Y \xrightarrow[\text{closed imm.}]{\ f^{\natural\prime}\ } C'(f) \xrightarrow{\ \lambda'_f\ } X,
  \qquad \sigma'_f : \text{open section and \textit{homotopism}}
\]
\note{comme à la page 10, les sections sont dessinées en arc de retour, et l'exposant de $f$ est lu $\natural$ sans certitude}

\note{en bas à droite, à l'encre}

discrete $\cap$ Serre ($=$ étale covering) $=$ smooth $\cap$ proper \struck{\ill{}} Serre
\note{au-dessus de « proper », au crayon, « étale » : le mot est ajouté puis laissé}

\struck{mono $\cap$ Serre} $=$ open $\cap$ closed imm.\ $=$ disjoint summand

\struck{discrete} $\cap$ \struck{\ill{}} Serre $=$ iso
\note{les premiers mots de ces deux lignes sont surchargés à l'encre ; la seconde se lit peut-être « discrete $\cap$ trivial Serre $=$ iso »}

\emph{NB} proper $\cap$ discrete $\Rightarrow$ cofib, smooth $\cap$ discrete $\Rightarrow$ fibration but does not imply mono, and proper \struck{\ill{}} Serre does not imply cofibration, \uncertain{nor trivial}, smooth $\cap$ Serre not imply fibration, \uncertain{nor trivial}
\note{comparer la note de la page 13}

\page{13}
\note{feuillet au crayon, en anglais ; deux colonnes séparées par un trait vertical, chacune divisée en trois cases par des traits horizontaux}

$f^{\sharp} \in$ fibration $\cap$ mono
\note{le signe en exposant, lu $\sharp$, pourrait être un $\natural$ ; il fait pendant au $f^{\flat}$ de la colonne de droite}

$\sigma_f \in$ cofibration $\cap$ \uncertain{monontr}

$\lambda_f \in$ not a fibration or cofibration nor Serre, but trivial fibers (with final object)

\[
  (Y \longrightarrow X) \Longrightarrow X \qquad Y_x \longrightarrow X \qquad 0 \longrightarrow 1
\]
\note{ces trois flèches, au crayon léger sous la case de $\lambda_f$, semblent illustrer le cas de $\lambda_f$ ; le second $X$ est repassé}

$(\lambda_f, \sigma_f)$ adjoint?

$(\mu_f, \rho_f)$ adjoint

$f^{\flat} \in$ fibration (not \emph{yet Serre})

$\rho_f \in$ cofibration $\cap$ Serre $\cap$ \struck{fib}

$\mu_f \in$ not a fibration nor cofibration nor Serre?

\marginal{NB proper discrete $\Rightarrow$ cofib ; \struck{\ill{}} smooth discrete $\Rightarrow$ fibration [but proper Serre $\not\Rightarrow$ cofib ; smooth Serre $\not\Rightarrow$ fib]}
\note{en haut à droite, séparé de la colonne par un trait courbe}

Dans une ellipse :
\[
  X \longrightarrow \text{cofibers of } f^{\sharp}\,? \qquad
  Y \longrightarrow \text{cofibers of } \sigma_f\,?
\]

\note{au-dessous, un tableau disposé en croix : la verticale, en lettres appuyées, va de \emph{bifibration} à \emph{mono} ; les deux bras portent le cas « fibration » à gauche et le cas « cofibration » à droite. Il est rendu ici ligne par ligne, à gauche puis à droite}

\[
  \begin{array}{c} \text{bifibration} \\ \uparrow \\ \text{Serre} \end{array}
\]

smooth \struck{\ill{}} fibration \quad $+$ \quad \uncertain{proper} $\Rightarrow$ co-proper
\note{à droite du $+$, plusieurs mots très effacés, où se devine deux fois \emph{proper}}

\[
  \begin{array}{c} \text{discrete} \\ \downarrow \\ \text{mono} \end{array}
\]

À gauche de \emph{discrete} : discrete fibration, $\hat{X}$ i.e. $X^{\circ} \to (\mathrm{Sets})$. À droite : discrete cofibration, $\check{X}$ i.e. $X \to (\mathrm{Sets})$.

À gauche de \emph{mono} :
\[
  X \xrightarrow{\ \varphi^{\circ}\ } [\emptyset, \{\emptyset\}] = \Delta_1, \qquad
  Y = \varphi^{-1}(\{1\}) = \varphi^{-1}(\{0\})
\]
\[
  (\psi : X \xrightarrow{\ \varphi^{\circ}\ } \Delta_1^{\circ} \simeq \Delta_1)
\]
open immersion
\note{au-dessus de $\emptyset$ et de $\{\emptyset\}$, les petits chiffres $0$ et $1$ ; la seconde égalité $\varphi^{-1}(\{1\}) = \varphi^{-1}(\{0\})$ est surchargée, la lecture de l'ordre des deux termes est incertaine}

À droite de \emph{mono} :
\[
  X \xrightarrow{\ \varphi'\ } [\emptyset, \{\emptyset\}] = \Delta_1 : [0 \longrightarrow 1], \qquad
  Y = \varphi'^{-1}(\{1\})
\]
closed immersion

\page{14}
Après l'histoire d'asphéricité, \ill{} compris en Loc~1, Loc~2, Loc~3 (Loc~1 sous forme faible :
\begin{itemize}
  \item[a)] iso dans $\underline{W}$
  \item[b)] si deux parmi $g$, $f$, $gf$ en $\underline{W}$, le 3\textsuperscript{e} aussi
  \item[c)] $gf = \mathrm{id}$ et $fg \in \underline{W} \Longrightarrow f, g \in \underline{W}$ )
\end{itemize}

\begin{itemize}
  \item[(0)] Axiomes Loc \uncertain{mettons en œuvre}
    \note{une flèche relie cette ligne à la fin de c) et à la ligne : « \ill{} \uncertain{transitions} fib et cofib. »}
  \item[(1)] \emph{Dualité} (y compris $\underline{W} = \underline{W}^{\circ}$)
  \item[(3)] Sorite des $\int \left(\begin{smallmatrix} Y & \to & X_0 \\ \downarrow & & \\ X_1 & & \end{smallmatrix}\right)$
\end{itemize}
\note{les numéros sont cerclés ; le 3 est repassé sur un 4}

\marginal{(2) Commutation aux fibrations, cofibrations, immersions ouvertes, fermées, cylindres, cônes \ldots}
\note{écrit en oblique dans la marge gauche, entouré d'un trait qui le rattache entre (1) et (3)}

\[
  \begin{array}{ccc}
    Y & \xrightarrow{\ f_0\ } & X_0 \\
    {\scriptstyle f_1}\downarrow & & \downarrow{\scriptstyle \alpha_1} \\
    X_1 & \xrightarrow{\ \alpha_0\ } & \Sigma
  \end{array}
  \qquad \Sigma \searrow X_0 \amalg_Y X_1
\]
\note{une diagonale biffée de $Y$ vers $\Sigma$, étiquetée $\lambda_g$ \uncertain{}; les étiquettes des flèches $\alpha_0$, $\alpha_1$ sont surchargées. À droite, traversé par un long trait courbe : « \uncertain{et propriétés} : \ill{} \uncertain{limite} »}

\struck{a) Si $f_0$ iso, \ill{} homotopismes}

\begin{itemize}
  \item[a)] \emph{Fonctorialité} si $u_Y, u_{X_0}, u_{X_1} \in \underline{W}$, alors $(\Sigma \to \Sigma') \in \underline{W}$
  \item[b)] Si $f_0$ iso, alors \struck{$\alpha_0$} $\alpha_0$ homotopismes (\ill{} $f_1$, $\alpha_1$)
  \item[c)] Si $f_0 \in \underline{W}$, alors $\alpha_0 \in \underline{W}$ (\ill{} $f_1$, $\alpha_1$)
\end{itemize}

\ill{} $f_0$, $f_1$ cofibrations (resp. fibrations). Supposons que $\forall x \in X$, $X_{0x} \to X_{1x}$ soit \struck{équival} $\underline{W}$-asphérique. Alors
\[
  \int \left(\begin{smallmatrix} Y & \to & X_0 \\ \downarrow & & \\ X_1 & & \end{smallmatrix}\right) \longrightarrow X
\]
\uncertain{est} cofib (resp. fib) : fibres triviales, donc $\in \underline{W}$.

\emph{Cor} Si $f_0 \in \underline{W}$, alors $q_0 \in \underline{W}$ ; si $f_1 \in \underline{W}$, alors $q_1 \in \underline{W}$.

(4) \struck{(3)} \emph{Thm} de \uncertain{« la somme amalgamée (cube) »}
\[
  \begin{array}{ccc}
    Y & \xrightarrow{\ f_0\ } & X_0 \\
    {\scriptstyle f_1}\downarrow & \text{cart} & \downarrow{\scriptstyle q_1} \\
    X_1 & \xrightarrow{\ q_0\ } & X
  \end{array}
\]
\note{le 4 cerclé est repassé sur un 3 ; l'étiquette $q_1$ est surchargée}

\emph{NB} Si \ill{} $\underline{W}$ \ill{} que $\forall x \in X$, \struck{$X_{0x} \to X_{1x}$ soit \ill{} cat. ponctuelle, \ill{}} \ill{} $X_{0x}$, $X_{1x}$ sont $\emptyset$, $\Delta_0$, l'autre est isom. : $\Delta_0$ \ill{} dans ce diagramme \ill{} cocartésien, et les \ill{} pour \ill{} simplifiés)
\note{passage très raturé, écrit en colonne étroite à gauche ; seuls quelques mots se lisent}

\emph{Question} Cas
\[
  \begin{array}{ccc}
    Y & \overset{i}{\hookrightarrow} & X \\
    {\scriptstyle f}\downarrow & & \downarrow{\scriptstyle g} \\
    Y' & \underset{i'}{\hookrightarrow} & X' = X \amalg_Y Y'
  \end{array}
  \qquad i \text{ immersion ouverte ou fermée}
\]
\uncertain{est-ce que} si $f$ cofib, il en résulte que $g$ le soit, et si $f \in \underline{W}$, que $g \in \underline{W}$ ? Il est possible qu'on n'ait pas $g \in \underline{W}$ [mais j'ignore si $i \in \underline{W}$ implique que $i' \in \underline{W}$]

(c) \struck{(4)} \emph{Application} aux sommes amalgamées dans les $A^{\wedge}$ (donc aux cartésiens et cocartésiens, avec $Y' \hookrightarrow X'$ mono)

4 \emph{Thm} Soit (de la somme \uncertain{amalgamée}, \uncertain{version} $A^{\wedge}$)
\[
  \begin{array}{ccc}
    Y' & \overset{i}{\hookrightarrow} & X \\
    {\scriptstyle g}\downarrow & & \downarrow{\scriptstyle f} \\
    Y' & \underset{i'}{\hookrightarrow} & X' = X \amalg_Y Y'
  \end{array}
\]
\note{au-dessus de la flèche du haut, « mono », souligné deux fois ; la source de la première ligne est écrite $Y'$, sans doute pour $Y$}
alors \struck{si $i \in \underline{W}$,}
\[
  i \in \underline{W}_A \Longrightarrow i' \in \underline{W}_A, \qquad
  g \in \underline{W}_A \Longrightarrow f \in \underline{W}_A
\]

(5) \emph{Application} : ($\mathrm{Hot}(\underline{W})$ et $\mathrm{Hot}_A$)

\emph{Thm} Calcul des \uncertain{fonctions} : \ill{} dans $\overrightarrow{\mathrm{Cat}}$, $\overrightarrow{A^{\wedge}}$
\note{le 5 cerclé est repassé sur un autre chiffre}

\page{15}
Application aux sommes et produits dans $\mathrm{Hot}(\underline{W})$, et les $\mathrm{Hot}_A$ ($A$ cat. test locale)

(6) \struck{\ill{}} Saturation \uncertain{forte} des $\underline{W}$, des $\underline{W}_A$.
\note{« forte » est ajouté au-dessus de la ligne ; la suite de la page est blanche}

\page{16}
\emph{Sorite des recollements}

$U \subset X \supset Y \qquad H_X : U^{\circ} \times Y \to \mathrm{Ens}$

a) Map \struck{\ill{}} $Z$ \quad $f : X \to Z$
\begin{itemize}
  \item[(i)] $f_U : U \to Z$ \quad \uncertain{induced by} $f$
  \item[(ii)] $f_Y : Y \to Z$ \quad \uncertain{induced by} $f$
  \item[(iii)] hom.\ of bifunctors
    \[ H_X(u, y) \longrightarrow \mathrm{Hom}_Z(f_U(u), f_Y(y)) \]
\end{itemize}

b) Map $f : X' \to X$, defining decomposition $X' = U'$ \struck{\ill{}} $\circ\, Y'$, with $H_{X'} : U'^{\circ} \times Y' \to (\mathrm{Ens})$, we must have:
\begin{itemize}
  \item[(i)] $f_{U'} : U' \to U$
  \item[(ii)] $f_{Y'} : Y' \to Y$
  \item[(iii)] bifunctorial $H_{X'}(u', y') \to H_X(u, y)$
\end{itemize}

c) Product
\[
  T \times U \subset T \times X \supset T \times Y \qquad
  H_{T \times X}((t, u), (t', y)) \simeq \mathrm{Hom}_T(t, t') \times H_X(u, y)
\]

d) homotopy $\Delta^1 \times X' \xrightarrow{\ h\ } X$ \emph{compatible} with decomposition
\note{le $\Delta$ est écrit sur une autre lettre, dans les deux facteurs ; le second facteur semble d'abord $X^1$, corrigé en $X'$}

i.e. $h^{-1}(U) = \Delta^1 \times U'$, $h^{-1}(Y') = \Delta \times Y'$ \quad $U', Y'$ \ill{}, $= h_i^{-1}(U)$, $= h_i^{-1}(Y)$, $i \in \{0, 1\}$
\note{les signes « $=$ » rendent deux guillemets de répétition placés sous $U'$ et $Y'$ ; on lit donc $U' = h_i^{-1}(U)$, $Y' = h_i^{-1}(Y)$. Le premier membre de la seconde égalité est écrit $h^{-1}(Y')$, pour $h^{-1}(Y)$ sans doute}

$h$ given by
\[
  h_{U'} : \Delta \times U' \to U, \qquad h_{Y'} : \Delta \times Y' \to Y
\]
plus \struck{bi} functorial map in $(s, u')$ $(t, y')$
\[
  \mathrm{Hom}(s, t) \times H_{X'}(u', y') \longrightarrow H_X(h_{U'}(s, u'), h_{Y'}(t, y'))
\]

\page{18}
\note{sa numérotation : « 1 » cerclé en tête de page ; la page 19 porte « 2 », la page 20 « 2 bis »}

\[
  \begin{array}{ccccc}
    Y & \overset{i}{\hookrightarrow} & X & \hookleftarrow & U \\
    {\scriptstyle g}\downarrow & & {\scriptstyle f}\downarrow & & \downarrow\wr \\
    Y' & \underset{i'}{\hookrightarrow} & X' & \hookleftarrow & U'
  \end{array}
\]

Soient $X'$ dans $(\mathrm{Cat})$, $Y'$ sous-objet fermé, $U'$ complémentaire ouvert, $f : X \to X'$, $Y = f^{-1}(Y')$, $U = f^{-1}(U')$. Supposons que $U \simeq U'$ ; alors \struck{$f$} le carré est-il néc.\ cocartésien ? \textbf{Non} bien sûr, comme on voit en prenant $X = U'$, $f$ l'inclusion, d'où $Y = \emptyset$.

\emph{Mais} supposons que $f$ soit cofibrant, le carré est-il cocartésien ? En d'autres termes, a-t-on pour $u' \in U'$, $y' \in Y'$
\[
  \mathrm{Hom}_{X'}(u', y') \overset{?}{\longleftarrow}
  \varinjlim_{y \in Y_{/y'}} \mathrm{Hom}(u, y)
\]
(où $u \in \mathrm{Ob}\, U$ au-dessus de $u'$). Sous la limite, une accolade renvoie à
\[
  \varinjlim_{Y_{y'}} \mathrm{Hom}(u, Y_{y'}) .
\]
\note{la flèche de la première formule est surmontée d'un « ? » et d'un $\sim$ ; entre les deux formules, un croquis : $u$ au-dessus de $u'$, $y$ au-dessus de $y'$ reliés par des pointillés verticaux, et une flèche pointillée de $u'$ vers $y'$ dont l'étiquette est raturée}

Oui, car \struck{$\mathrm{Hom}(u,$} (i) surjectif et (ii) pour \struck{$y$} : $u \to y'$ fixé \ldots

Nous nous intéressons donc à la situation des diagrammes du type précédent, avec $f$ \emph{cocartésien}. Pour se donner un foncteur $X' \to (\mathrm{Cat})$, \struck{donc} dont la restriction à $U'$ soit le foncteur c\textsuperscript{t} de valeur $\Delta_0 = e_{(\mathrm{Cat})}$, i.e.
\begin{itemize}
  \item[(i)] $Y' \xrightarrow{\ \Phi\ } \mathrm{Cat}$ (définissant \struck{f} $Y \to Y'$)
  \item[(ii)] $\mathrm{Hom}(u', y') \to \mathrm{Ob}\, \Phi(y')$ bifonctoriel en $u', y'$ en $U' \times Y'$
\end{itemize}

\page{19}
On voudrait que sous ces conditions
\[
  \int \left(\begin{smallmatrix} Y & \to & X \\ \downarrow & & \\ Y' & & \end{smallmatrix}\right)
  \longrightarrow X' = \varinjlim \left(\begin{smallmatrix} Y & \to & X \\ \downarrow & & \\ Y' & & \end{smallmatrix}\right) = X \amalg_Y Y'
\]
soit $\in \underline{W}$, et \add{même} qu'il soit asphérique.
\note{« même » est rendu par son signe $=$ ; le passage est marqué d'un double trait vertical dans la marge}

Les conditions sont \struck{elles} stables par \emph{tout} changement de base sur $X'$.

\[
  \begin{array}{ccc}
    \int(i, g) & \longrightarrow & X' \\
    \uparrow & & \uparrow \\
    \int(i_1, g_1) & \longrightarrow & X'_1
  \end{array}
\]
\note{à côté, deux croquis d'un cube de changement de base : la face $Y \overset{i}{\hookrightarrow} X$, $g$, $f$, $Y' \overset{i'}{\hookrightarrow} X'$, et la face $Y_1 \overset{i_1}{\hookrightarrow} X_1$, $g_1$, $f_1$, $Y'_1 \overset{i'_1}{\hookrightarrow} X'_1$, reliées par des flèches obliques ; un troisième croquis, à gauche, est raturé}

Ce diagramme est-il cartésien ? À vérifier ! Si oui, cela signifierait que $\int(i, g) \to X'$ est \emph{universellement} dans $\underline{W}$.

\emph{Mais} on vérifie directement que $\int(i, g) \to X'$ est \emph{cofibré} à fibres $\underline{W}$-asphériques ; en fait, fibres sont
\[
  \simeq e_{u'} \ \text{pour } u' \in U', \qquad
  \simeq C(Y_{y'}) \ \text{pour } x = y' \in Y'
\]
\note{devant $u' \in U'$, « $x =$ » est biffé}
(\uncertain{cône}, \ill{} \uncertain{élément final} : $Y_{y'}$ !)
\note{ce paragraphe est marqué d'un triple trait vertical dans la marge}

\page{20}
\note{« 2 bis » en tête de page}

\marginal{\uncertain{renvoyer} aux fibrations}
\note{souligné, en oblique en haut de la marge gauche}

\emph{Thm} Soit
\[
  \begin{array}{ccc}
    Y & \xrightarrow{\ i\ } & X \\
    {\scriptstyle g}\downarrow & & \downarrow{\scriptstyle f} \\
    Y' & \xrightarrow{\ i'\ } & X'
  \end{array}
\]
un diagramme cocartésien (dans $(\mathrm{Cat})$), avec $i'$, $f$ des cofibrations, \struck{et} supposons que $\forall x' \in X'$, $f^{-1}(x')$ \emph{ou} $i'^{-1}(x')$ soit équiv.\ à la catégorie ponctuelle, et considérons le foncteur canonique
\[
  \int \left(\begin{smallmatrix} Y & \to & X \\ \downarrow & & \\ Y' & & \end{smallmatrix}\right) \longrightarrow X' .
\]
\marginal{au lieu de \uncertain{ponctuelle}, il suffit \ill{} soit asphérique}
\note{« (dans (Cat)) » est ajouté au-dessus de la ligne ; dans l'intégrale, un $X'$ en bas à droite est biffé}

Ce foncteur est une cofibration, à fibres les
\[
  \int \left(\begin{smallmatrix} Y_{x'} = Y'_{x'} \times X_{x'} & \xrightarrow{\ \mathrm{pr}_2\ } & X_{x'} \\ \downarrow{\scriptstyle \mathrm{pr}_1} & & \\ Y'_{x'} & & \end{smallmatrix}\right) ,
\]
qui sont \struck{\ill{}} de la forme
\[
  \int \left(\begin{smallmatrix} A & \xrightarrow{\mathrm{id}_A} & A \\ \downarrow & & \\ e & & \end{smallmatrix}\right)
  \qquad \left(\text{ou } \int \left(\begin{smallmatrix} A & \to & e \\ \downarrow{\scriptstyle \mathrm{id}_A} & & \\ A & & \end{smallmatrix}\right)\right),
\]
donc asphériques.
\note{au-dessus de « de la forme », un ajout raturé, où se lit « des cat. » ; dans la première intégrale, l'égalité $Y_{x'} = Y'_{x'} \times X_{x'}$ est surchargée}

La catégorie \ill{} $C(A) \amalg_A (A \times \Delta^1)$ \struck{il y a une rétraction \ill{}} $C(A)$ en est un rétracte par déformation, on gagne !

\struck{Il suffit de voir que} $\int (A \to A)$ et $\int (A \to e)$ sont $\underline{W}$-asphériques (par propriété des recollements), \ill{} $Z$. \ill{} fini, \ill{} c'est clair \ill{} pour le premier \ill{} des $\underline{W}$ \ill{} $\underline{W}$
\note{ce passage est enclos d'un long trait et barré de deux traits obliques ; la lecture en est très fragmentaire}

\emph{Lemme} Si dans $\int \left(\begin{smallmatrix} Y & \xrightarrow{i} & X \\ {\scriptstyle g}\downarrow & & \\ Y' & & \end{smallmatrix}\right)$ \ill{}
\note{écrit en oblique au pied de la page, inachevé}

\emph{Corollaire}
\[
  i \in \underline{W} \Longrightarrow i' \in \underline{W}, \qquad
  g \in \underline{W} \Longrightarrow f \in \underline{W}
\]

\end{document}
